%!TEX root=../../assign.tex
\subsection{Problem 4}
\subsubsection{\rom{1}}
\begin{proof}
    Suppose 
    \begin{equation*}
        \paren{\forall x > 0 }\paren{\exists n \in \N } \frac{1}{n} < x. 
    \end{equation*}

    Multiplying both sides by $\frac{n}{x}$, 
    \begin{equation*}
        \frac{1}{x} < n. 
    \end{equation*}

    Since every positive real number has a positive multiplicative inverse, we have
    \begin{equation*}
        \paren{\forall x > 0}\paren{\exists x' > 0} x \cdot x' = 1. \quad \implies \quad x' = \frac{1}{x} > 0
    \end{equation*}

    This yields 
    \begin{equation*}
        \paren{\forall x > 0}\paren{\exists x' > 0}\paren{\exists n \in \N} \frac{1}{x} = x' < n, 
    \end{equation*}

    Now, Since we also have 
    \begin{equation*}
        \paren{\forall x' > 0 }\paren{\exists x > 0 }x' = \frac{1}{x} > 0, 
    \end{equation*}

    this yields
    \begin{equation*}
        \paren{\forall x' > 0 }\paren{\exists n \in \N} x' < n. 
    \end{equation*}

    Thus, the Archimedean property is proved. 
\end{proof}

\newpage
\subsubsection{\it{ii}}
\textbf{Proof for the case for $x \in \Q$}

We need to prove 
\begin{equation*}
    \paren{\forall x \in \Q}\bracks{\paren{x > 0}\implies \paren{\exists n \in \N}\frac{1}{n} < x}. 
\end{equation*}
\begin{proof}
    Since $x \in \Q$, by the definition
    \begin{equation*}
        \paren{\forall x \in \Q}(\exists m, d \in \Z) \bracks{d \neq 0 \implies x = \frac{m }{d}}. 
    \end{equation*}

    Since $x > 0$, which implies $m \neq 0$. 

    Consider the case where both $m$ and $d$ are positive first. We have 
    \begin{equation*}
        \frac{1}{n} < \frac{m }{d} \quad \implies \quad d < mn. 
    \end{equation*}

    Since $d \geq 1$ and $m \geq 1$, we can set 
    \begin{equation*}
        n = d + 1. 
    \end{equation*}

    This yields, 
    \begin{equation*}
        d \leq md < m(d + 1) = md + m. 
    \end{equation*}

    Therefore, 
    \begin{equation*}
        d < mn \quad \implies \quad \frac{1}{n} < \frac{m }{d}. 
    \end{equation*}

    Thus, the original claim holds. 

    Consider the case where $m, d < 0$. 

    Since 
    \begin{equation*}
        \frac{m  }{d} = \frac{-m }{-d}, 
    \end{equation*}
    we are able to make 
    \begin{equation*}
        m' = -m \quad \text{and} \quad d' = -d
    \end{equation*}

    Then treat them as the case where $m$ and $d$ are both positive. 
\end{proof}

\textbf{Prove for the case where $x > 0, x = \sqrt{y}, \exists y \in \Q$}. 
We need to prove 
\begin{equation*}
    \paren{\forall x > 0}\bracks{\bracks{\paren{\exists y \in \Q}\bracks{y > 0 \wedge x = \sqrt{y}}} \implies \paren{\exists n \in \N } \frac{1}{n} < x}.
\end{equation*}
\begin{proof}
    Let 
    \begin{equation*}
        y = \frac{m }{d}, \exists m \in \Z, \exists d \in \Z, d \neq 0. 
    \end{equation*}

    Since $y > 0$, which implies 
    \begin{equation*}
        m \neq 0, 
    \end{equation*}

    and we have two cases to discuss. They are 
    \begin{equation*}
        \paren{m \geq 1}\wedge\paren{d \geq 1} \quad \text{and} \quad \paren{m \leq -1}\wedge\paren{d \leq -1}. 
    \end{equation*}

    \begin{itemize}
        \item Case $\paren{m \geq 1}\wedge\paren{d \geq 1}$: 
        \begin{align*}
            \paren{\exists n \in \N}\frac{1}{n} < x &\Leftrightarrow \paren{\exists n \in \N}\frac{1}{n} < \sqrt{y} = \sqrt{\frac{m }{d}} \\
            &\Leftrightarrow \paren{n \in \N } \frac{1}{n^2} < \frac{m }{d} \\
            &\Leftrightarrow \paren{\exists n \in \N } d < mn^2,  
        \end{align*}

        Since $m \geq 1$, we can select $n = \paren{d + 1}$. We have 
        \begin{equation*}
            mn^2 = m\paren{d + 1}^2 = m\paren{d^2 + 2d + 1} = md^2 + 2md + m > d
        \end{equation*}

        Thus, the claim hold. 
        \item Case $\paren{m \leq -1}\wedge \paren{d \leq -1}$
        
        Since $y > 0$, we are able to multiple $y$ by $\frac{-1}{-1} = 1$, where we have 
        \begin{equation*}
            y = \frac{m }{d} \cdot 1 = \frac{m }{d} \cdot \frac{-1}{-1} = \frac{-m }{-d}, -m \geq 1, -d \geq 1. 
        \end{equation*}

        Then, we treat this case similar to the last case. 
    \end{itemize}

    Thus, the claim holds. 
\end{proof}

\newpage
\subsubsection{\rom{3}}
\begin{proof}
    Suppose $x$ is an arbitrary positive number in $\R$. We define the set 
    \begin{equation*}
        T := \set{m \in \N : x < m}. 
    \end{equation*}

    By Archimedean Property, we know that $T \neq \emptyset$. By well ordering principle, since $T \neq \emptyset$ and $T \subset \N$, $T$ has a least element. We define 
    \begin{equation*}
        k := \min T. 
    \end{equation*}

    Since $k$ is the least element in $T$, which means 
    \begin{equation*}
        k > x \quad \text{and} \quad k - 1 \not\in T. 
    \end{equation*}

    The definition of set $T$ implies 
    \begin{equation*}
        k - 1 \leq x. 
    \end{equation*}

    Thus, we have 
    \begin{equation*}
        k - 1 \leq x < k. 
    \end{equation*}

    Let $n = k - 1$, we replace $k$ by $n$. 
    \begin{equation*}
        n \leq x < n + 1. 
    \end{equation*}

    Thus, the original claim holds. 

    For $x < 0$ in $\R$. You are able to add a natural number $t > -x$ which makes $t + x$ positive. 
    Following the similar steps as $x > 0 \in \R$. We prove there exists a number $n$ in $\N$ such that 
    \begin{equation*}
        n \leq t + x < n + 1. 
    \end{equation*}

    Subtracting $t$, we get 
    \begin{equation*}
        n - t \leq x < n - t + 1. 
    \end{equation*}

    Let us define $z = n - t$, we get 
    \begin{equation*}
        z \leq x < z + 1
    \end{equation*}

    where $z \in \Z$. 

\end{proof}

At here, we prove the uniqueness of such integer $m$. 
\begin{proof}
    Suppose there exists two integers $n$ and $m$ in $\N$ such that 
    \begin{equation*}
        n \leq x < n + 1 \quad m \leq x < m + 1
    \end{equation*}
    and 
    \begin{equation*}
        n \neq m. 
    \end{equation*}

    Without loss of generality, we assume $n < m$. We have 
    \begin{equation*}
        n + 1 \leq m. 
    \end{equation*}

    Then, we get 
    \begin{equation*}
        n \leq x < n + 1 \leq m \leq x < m + 1. 
    \end{equation*}

    This yields, 
    \begin{equation*}
        x < x,  
    \end{equation*}

    which contradicts the fact that $x = x$. 

    Thus, $n = m$. 
    Therefore, the integer $m$ is unique. 
\end{proof}

\newpage
\subsubsection{\rom{4}}
\begin{proof}
    Suppose $x$ and $y$ are both arbitrary number in $\R$ with $y > x \geq 0$. 
    
    By hypothesis, 
    \begin{equation*}
        y - x > 0. 
    \end{equation*}

    Following Archimedean Property, there exists $n \in \N$ such that 
    \begin{equation*}
        y - x > \frac{1}{n} > 0. 
    \end{equation*}

    This yields
    \begin{equation*}
        ny - nx > 1,  
    \end{equation*}

    which implies 
    \begin{equation*}
        ny > 1 + nx. 
    \end{equation*}

    Moreover, by Archimedean Property, there exists $m \in \N$ such that  
    \begin{equation*}
        m - 1 \leq nx < m. 
    \end{equation*}

    Adding $1$, we obtain
    \begin{equation*}
        m \leq nx + 1 < m + 1. 
    \end{equation*}

    Hence 
    \begin{equation*}
        m - 1 \leq nx < m  \leq nx + 1 < ny, 
    \end{equation*}

    which implies 
    \begin{equation*}
        nx < m < ny. 
    \end{equation*}

    Dividing $n$, 
    \begin{equation*}
        x < \frac{m }{n} < y. 
    \end{equation*}

    By definition, 
    \begin{equation*}
        x < q < y   
    \end{equation*}

    where $q = \frac{m }{n } \in \Q$. 

    If
    \begin{equation*}
        x < 0, 
    \end{equation*}
    then there exists a natural number $k$ with $k > -x$ such that 
    \begin{equation*}
        \paren{\exists m, n \in \N} x + k < \frac{m }{n} < y + k, 
    \end{equation*}

    where $x + k > 0$ and $y + k > 0$. 

    Subtracting $k$ 
    \begin{equation*}
        x < \frac{m }{n} - k < y. 
    \end{equation*}

    We have 
    \begin{equation*}
        x < \frac{m - nk}{n} < y 
    \end{equation*}
    
    By definition, there exists $q$ such that 
    \begin{equation*}
        x < \frac{m - nk}{n} = q < y 
    \end{equation*}

    where $m - nk \in \Z$ and $n \in \N$.
    
    This proves the claim. 
\end{proof}